\documentclass[11pt]{article}
\usepackage[margin=1in]{geometry}
\usepackage{amsmath}
\usepackage[version=4]{mhchem}
\usepackage[T1]{fontenc}
\usepackage{lmodern}
\usepackage[authoryear,round]{natbib}
\usepackage{graphicx}
\usepackage{caption}
\usepackage{url}
\usepackage[left,mathlines]{lineno}
\usepackage[hidelinks]{hyperref}
\setlength{\emergencystretch}{2em}

\title{Precession-phase preference of Dansgaard--Oeschger event occurrence after accounting for background climate}
\author{}
\date{}
\begin{document}
\maketitle
\linenumbers

\section*{Key Points}
\begin{itemize}
\item Precession phase improves warming-onset fits after accounting for background climate and recent events.
\item The association generally persists under age perturbations, but finite event counts leave broad uncertainty in effect strength and phase.
\item Overlapping fit improvements from phase and summer insolation, together with age-model dependencies, limit identification of a unique orbital mechanism.
\end{itemize}

\begin{abstract}
Orbital forcing modulates millennial climate variability, but its relationship to individual abrupt transitions remains uncertain. We analyze 55 Greenland and speleothem warming onsets, using 53 response events to fit a continuous point process conditional on background climate and exponentially decaying event history. Adding precession phase gives a log-likelihood gain per event of 0.248 bits/event, with a reduced-model bootstrap $p=0.0003$ and maximum occurrence near the precession-index minimum. The association generally persists under chronology perturbations, although event simulations reveal broad uncertainty in effect strength and preferred phase. A synthetic Greenland catalogue gives a similar phase preference across 69 response events (bootstrap $p=0.0077$). Phase and summer-solstice insolation produce overlapping fit improvements. These results support phase-dependent warming occurrence, while finite event counts and chronology dependencies limit attribution to a unique orbital mechanism.
\end{abstract}

\section*{Plain Language Summary}
Ice cores and cave deposits record repeated abrupt climate changes during past ice ages. Changes in Earth's orbit may influence when these transitions occur. We examine 55 warming events, allowing their occurrence rate to depend on background climate and earlier events. The model better describes their timing when it includes precession, the slow change in the orientation of Earth's axis relative to its orbit. This relationship usually persists when event ages are varied within the assumed uncertainty, and a separate reconstruction shows a similar pattern. Some orbital configurations therefore appear to favor abrupt warming. However, the few available events leave uncertainty in the strength and preferred timing of this effect. The results also cannot distinguish precession from the summer sunlight changes it produces, or identify a unique physical trigger.

\section{Introduction}

Dansgaard--Oeschger (DO) events punctuated the last glacial period with abrupt Greenland warmings from cold stadials into warmer interstadials, followed by more gradual cooling \citep{dansgaard1993evidence}. They are a prominent expression of millennial-scale climate variability (MCV). Using synchronized oxygen-isotope and calcium records from three Greenland ice cores, \citet{rasmussen2014stratigraphic} established a consistent event stratigraphy, defining stadial and interstadial boundaries on a common chronology. This framework provides a reference sequence for investigating abrupt warming occurrence during the last glacial period.

Speleothems extend the evidence for DO-like variability beyond this Greenland reference sequence. Uranium--thorium (U--Th) chronologies show that last-glacial Greenland warmings were accompanied by near-synchronous hydroclimate changes across monsoon and European--Mediterranean regions \citep{corrick2020synchronous}. Older records from Melchsee--Frutt in the Swiss Alps and Sofular Cave in T\"urkiye document DO-like stadial--interstadial transitions during the penultimate glacial \citep{fohlmeister2023role,held2024dansgaard}. These regional records complement Greenland ice cores, providing dated event sequences for comparing abrupt climate variability across successive glacials.

Intermediate glacial conditions are particularly favorable for the MCV \citep{barker2011800,malmierca2023dansgaard}. Climate-model experiments show that these background conditions can support oscillations of the Atlantic Meridional Overturning Circulation (AMOC) between stronger and weaker states, producing abrupt stadial--interstadial transitions through various mechanisms \citep{vettoretti2022atmospheric,malmierca2023dansgaard}.

Orbital modulation is evident in the magnitude of MCV. After orbital-scale fluctuations are removed from proxy records, the amplitude envelopes and moving-window standard deviation of the remaining millennial variability can still contain orbital periodicities. A synthesis spanning 1.5 Myr identifies both obliquity- and precession-band modulation, alongside stronger glacial modulation after the mid-Pleistocene transition \citep{sun2021persistent}. These relationships differ among archives: \citet{thirumalai2020methane} found pronounced precession modulation of millennial methane variability but no corresponding precession signal in the millennial amplitude envelope of a composite Chinese speleothem record. The orbital imprint thus appears to depend on the regional climate processes and proxy responses represented by each archive.

Amplitude modulation, however, does not by itself establish orbital regulation of event occurrence. Such modulation could arise if orbital forcing changed event magnitudes while leaving the probability of occurrence unchanged. Climate-model experiments nevertheless provide a physical basis for testing an occurrence response: changing orbital parameters can initiate DO-like AMOC oscillations under fixed ice-sheet and greenhouse-gas conditions, through changes in tropical hydroclimate and subpolar sea ice \citep{zhang2021direct}. These oscillations can persist under constant orbital settings. The experiments therefore demonstrate orbital influence on the conditions permitting abrupt variability, while leaving its relationship to the timing of observed warmings to be tested.

Previous statistical and modeling studies have already linked transition occurrence to climate and insolation. For last-glacial Greenland, \citet{lohmann2018random} found that warming counts alone were compatible with a constant-rate Poisson process, whereas an alternating model with constant warming and cooling rates could not reproduce the observed fluctuations. Their forced alternating model associated higher boreal summer insolation with shorter stadials and greater ice volume with shorter interstadials. Across the last two glacials, \citet{fohlmeister2023role} identified interstadial-onset clusters near summer-insolation maxima and improved simulations of transition counts and warm-state occupancy by combining sea-level forcing with an insolation-dependent warming threshold. These findings motivate a more specific statistical question: do orbital parameters add information about individual warming onsets after accounting for background climate and recent events? Establishing this requires a conditional comparison that also quantifies the constraints imposed by uncertain ages and limited event counts.

Precession contributes to summer-insolation variability by changing the timing of the seasons relative to perihelion \citep{zhang2021direct}, motivating our focus on its phase. Here we combine the last-glacial Greenland sequence with a composite speleothem event sequence centered on Marine Isotope Stage (MIS) 6. The synthetic Greenland reconstruction of \citet{barker2011800}, on its speleothem-based chronology, provides a separate comparison. We fit a continuous-time point process to warming onsets, using the LR04 benthic $\delta^{18}$O stack and atmospheric \(\ce{CO2}\) as indicators of background climate, together with exponentially decaying event history. Nested models with and without precession phase test its additional contribution to occurrence per unit of total observed time. We estimate the preferred phase and modulation strength, assess uncertainty from event ages and finite event counts, and compare phase with other orbital predictors. This approach tests whether the event sequence carries a phase association beyond background climate, complementing evidence from variability amplitudes and climate-state durations.

\section{Data and methods}
Our primary sequence combines 34 last-glacial interstadial starts from the North Greenland Ice Core Project (NGRIP; \citealp{rasmussen2014stratigraphic}) with a 21-event speleothem sequence centered on MIS~6 and spanning approximately 135--194 kyr before 1950 CE (kyr BP; Figure~\ref{fig:data_framework}b). Sixteen speleothem transitions come from the Melchsee--Frutt (MF) oxygen-isotope stack in the Swiss Alps \citep{fohlmeister2023role}, and five form an older continuation from the Sofular carbon-isotope record in T\"urkiye \citep{held2024dansgaard}. Their high resolution and corresponding millennial-scale features across the overlapping interval allow distinct events to be combined without adjusting either chronology (Figure~S1). We retain published event identities and assign transition times from the strongest proxy gradients near the original labels (Text~S1; Figure~S2; Table~S1). 

The Greenland and speleothem sequences are fitted jointly, without carrying event history across the gap between them. Fitting starts immediately after the oldest event in each sequence; that event initializes the history term but is not counted among the fitted events. This leaves 53 of the 55 events for rate estimation. The interval between the youngest recorded warming and the younger boundary of each record remains included in the fit (Text\~S1).

For a complementary test over a longer interval, we use the Antarctic-derived synthetic Greenland reconstruction of \citet{barker2011800}. Its published catalogue contains 70 variable-threshold warming onsets over 0--400 kyr BP on SpeleoAge (a chronology aligned to independently U--Th-dated Chinese speleothem records). These events form the main Barker analysis; a 59-event fixed-threshold subset tests sensitivity to event definition (Figure~\ref{fig:data_framework}b; Text~S1).

\begin{figure}[p]
\centering
\includegraphics[width=\linewidth,height=0.68\textheight,keepaspectratio]{figures/Fig01.pdf}
\caption{Event inventories and environmental context. (a) Eccentricity, obliquity, and daily mean summer-solstice insolation at $65^\circ$N from La2004 \citep{laskar2004long}; axes match curve colors. (b) Precession index with NGRIP--MIS6 events (blue triangles; 55), Barker varying-threshold events (rose circles; 70), and the fixed-threshold subset (open green squares; 59). Shared Barker events have overlapping symbols. Blue strips mark the observation intervals of NGRIP--MIS6 events; Barker spans 0--400 kyr BP. (c) LR04 benthic $\delta^{18}$O \citep{lisiecki2005pliocene}, with larger values downward. (d) Atmospheric CO$_2$ \citep{bereiter2015revision}. Event sources are \citet{rasmussen2014stratigraphic,fohlmeister2023role,held2024dansgaard,barker2011800}.}
\label{fig:data_framework}
\end{figure}

We represent background climate using LR04 benthic $\delta^{18}$O and the Antarctic atmospheric CO$_2$ composite \citep{lisiecki2005pliocene,bereiter2015revision}. LR04 reflects both ice volume and deep-water temperature. Orbital inputs follow La2004 \citep{laskar2004long}. Forcing ages remain fixed during event-age perturbations for sensitivity analyses of chronology uncertainty (Text S2). Phase is interpolated between alternating precession-index extrema: minima define $0^\circ$, maxima $180^\circ$, and phase increases toward older ages.

We fit a continuous point process whose conditional intensity combines climate and event history \citep{truccolo2005point}. Let $u$ denote time in kyr advancing from a segment's conditioning event toward the present. The instantaneous warming rate is
\begin{equation}
\log\lambda_r(u)=\beta_0+\beta_H H_r(u)+\beta_L L_r(u)+\beta_C C_r(u)+\beta_S S_r
+\beta_s\sin\phi_r(u)+\beta_c\cos\phi_r(u),\label{eq:rate}
\end{equation}
where $H_r(u)=\sum_{j:0\leq u_{rj}<u}\exp[-(u-u_{rj})/\tau]$ includes strictly earlier events of the same segment, including its conditioning event at $u=0$. The fixed decay time is $\tau=1.5$ kyr. We constrain $\beta_H\leq0$, allowing recent events to suppress the rate, and set unobserved pre-anchor history to zero. $L_r$ and $C_r$ denote LR04 and CO$_2$, and $S_r$ distinguishes the primary segments; Barker omits $S_r$. Continuous interpolants give external covariates at event and integration times. Climate scaling uses nominal exposure-time means and ranges.

We estimate coefficients by maximizing the conditional event likelihood,
\begin{equation}
\ell=\sum_r\left[\sum_{i=1}^{N_r}\log\lambda_r(u_{ri})-\int_0^{T_r}\lambda_r(u)\,du\right],\label{eq:likelihood}
\end{equation}
where $T_r$ is response exposure and event intensities use the history immediately before each event. Numerical quadrature resolves actual event ages and external interpolation knots (Text~S3). The conditioning event contributes to history but not the event sum. This formulation represents warming onsets per unit of total response time.

The background model (BG) omits the two phase terms. Both models are refitted, and their likelihood improvement is expressed as
\begin{equation}
G=\frac{\ell_{\mathrm{full}}-\ell_{\mathrm{BG}}}{N\ln2},\qquad
LR=2(\ell_{\mathrm{full}}-\ell_{\mathrm{BG}}),\label{eq:gain}
\end{equation}
where $N$ is the response event count. The log-likelihood gain per event, $G$, measures in-sample fit improvement in bits/event; it does not establish held-out predictive performance. Positive $G$ alone is insufficient evidence because additional parameters can improve nested-model fits under the null. The fitted phase maximum is $\operatorname{atan2}(\beta_s,\beta_c)$ and the maximum/minimum phase rate ratio is $\exp[2(\beta_s^2+\beta_c^2)^{1/2}]$. Rayleigh statistics describe unadjusted event-phase concentration separately (Text~S3).

We calibrate each main association using 9,999 parametric bootstrap simulations under its fitted BG model \citep{davison1997bootstrap}. Each simulation retains the exact conditioning events and response support, generates continuous event times with immediate history updates, and refits both models. Event totals vary. Empirical $p=(1+\#\{LR^*\geq LR\})/10{,}000$ calibrates the added phase block; other likelihood-ratio tests use nominal asymptotic probabilities unless stated otherwise. Separate simulations test the history contribution and assess full-model fit using rescaled event intervals (Text~S8).

Chronology Monte Carlo (MC) ensembles preserve event identities and order (Texts~S2--S4). NGRIP uses cumulative Gaussian age displacements constrained by layer-counted errors, together with event-definition errors \citep{rasmussen2022gicc05,svensson2008sixty}. Beyond the counted section, published uncertainty is unquantified \citep{corrick2020synchronous}. For the modeled section, we adopt a uniform relative envelope of $\pm4.5\%$ of b2k age, interpreted approximately as $2\sigma$, based on the uncertainty estimate at 120 ka in \citet{moseley2020nalps19}. MF uses a shared age-error factor; Sofular uses dated-depth perturbations transferred through individual age models. Nine smoothing/window combinations represent speleothem event-definition sensitivity. Barker uses uniform proposals within published combined control uncertainties, interpolated on SpeleoAge. Each main catalogue has 10,000 age realizations, with unsupported draws flagged rather than replaced. Each valid draw has its own exact oldest conditioning event and response exposure, while retaining nominal climate scaling. We refit at the realized event ages and rebuild their history.

For primary-catalogue effect precision, we compare three constructions (Text~S5): age-only ranges (A); an approximate joint coefficient confidence region calibrated by 5,000 full-model simulations at nominal ages (B); and working ranges from 50 full-model simulations for each of 200 valid chronology generators (C). C represents combined sensitivity, not a calibrated combined confidence interval. Sensitivity tests vary history decay and initialization, background-response shape, event definition and segment phase response (Text~S6). Additional random deletion of 10\% or 20\% of noninitial events assesses sensitivity to further inventory loss (Text~S8). Orbital comparisons add eccentricity, obliquity or 65$^\circ$N summer-solstice insolation, with and without phase. The nine additional comparisons per catalogue form a separate nominal Holm family \citep{holm1979simple}. LR04--phase interactions are exploratory (Text~S7).

\section{Results}
Precession phase improves the primary event-rate fit by $G=0.2482$ bits per response event ($LR=18.238$), with bootstrap $p=0.0003$ and nominal $p=0.000110$ (Figure~\ref{fig:phase_comparison}; Figure~S8a). The maximum occurs at $328.1^\circ$, near the precession-index minimum, with a maximum/minimum rate ratio of 7.18 at fixed background and history. The unadjusted 55-event Rayleigh test gives $p=0.058$. Unlike that descriptive test, the conditional comparison accounts for climate, history and exposure using 53 response events. Overlaying the fitted full-model curve on the raw counts shows it bulging toward the observed peak near the preferred phase, while the no-phase model stays nearly flat (Figure~\ref{fig:phase_comparison}a,c).

Barker variable-threshold events favor a similar phase, with $G=0.1106$ bits per response event, bootstrap $p=0.0077$, a maximum at $334.2^\circ$ and a smaller rate ratio of 3.32 (Figure~S8b). Fixed-threshold picks give $G=0.1067$, nominal $p=0.0137$, a maximum at $347.1^\circ$ and a rate ratio of 3.19 (Figure~\ref{fig:phase_comparison}c,d). Thus the preferred phase sector persists across event definitions. Their 69 and 58 response events use different conditioning anchors and exposure, so the $G$ comparison is not a same-support likelihood ranking. The fixed-threshold result is a point-age sensitivity without additional bootstrap calibration.

\begin{figure}[p]
\centering
\includegraphics[width=\linewidth,height=0.775\textheight,keepaspectratio]{figures/Fig02.pdf}
\captionsetup{font=footnotesize}
\caption{Precession-phase organization and conditional phase response. Blue, rose, and green identify NGRIP--MIS6, Barker varying threshold, and Barker fixed threshold. (a,c) Unadjusted event counts in $20^\circ$ sectors (bars), overlaid with fitted expected counts per sector from the no-phase (dashed) and full (solid) models for the primary and Barker varying-threshold catalogues. The full model's curve bulges toward the observed counts near the preferred phase, while the no-phase curve stays nearly flat. (b,d) Conditional phase multipliers $\exp(\beta_s\sin\phi+\beta_c\cos\phi)$, with maxima marked by dots; ticks along the bottom mark inventory-event phases. Annotations give log-likelihood gain per event, $G$ (bits/event), nominal likelihood-ratio $p$, preferred phase, and maximum/minimum rate ratio. Unity represents zero phase contribution, rather than the refitted reduced-model rate. Minimum/Maximum labels identify precession-index extrema, not rate extrema. Histograms include 55/70/59 inventory events, whereas fitted curves and conditional fits use 53/69/58 response events. All results use nominal ages; sampling and chronology uncertainty are assessed separately.}
\label{fig:phase_comparison}
\end{figure}

Of 10,000 primary age realizations, 9,982 remain within observation support; 18 are excluded. Median $G$ is 0.2247 bits per response event, with a central 95\% chronology-sensitivity range of 0.1241--0.2914. Nominal $p<0.05$ in 9,961/9,982 valid fits (99.79\%; Figure~S6). All 10,000 Barker fits are valid, with median $G$ 0.0986 (95\% range 0.0708--0.1308) and a corresponding fraction of 99.48\% (Figure~S7). These fractions measure sensitivity to the adopted age perturbations, not bootstrap significance.

Full-model simulations show broader effect uncertainty than chronology perturbations alone (Figure~\ref{fig:effect_precision}). Primary age-only ranges are $311.6^\circ$--$342.8^\circ$ for preferred phase and 3.72--9.20 for the rate ratio. The sampling confidence-region projections span $287.8^\circ$--$370.2^\circ$ and 1.73--29.88; angles above $360^\circ$ continue into the next cycle. Combined working ranges are $297.0^\circ$--$355.8^\circ$ and 2.80--27.37. The phase preference remains broadly localized, but the magnitude of modulation is imprecise. B and C need not be nested because they respectively project a joint confidence region and summarize marginal refit ranges.

\begin{figure}[p]
\centering
\includegraphics[width=\linewidth,height=0.68\textheight,keepaspectratio]{figures/Fig03.pdf}
\caption{Uncertainty in the primary precession effect. A: central 95\% ranges across 9,982 chronology fits. B: projections of an approximate joint 95\% coefficient confidence region calibrated with 5,000 full-model simulations at nominal ages. C: central 95\% working ranges from 50 simulations for each of 200 chronology generators (Text~S5). (a,b) Preferred phase, unwrapped around $328.1^\circ$, and maximum/minimum rate ratio. Dashed lines show nominal estimates; dots show A/C medians and the B point estimate. (c) B coefficient-region boundary, every fifth C refit, nominal estimate (black dot), and zero effect (cross). (d) Nominal phase multiplier (black), simultaneous B confidence envelope (blue shading), and pointwise C range (orange dotted lines). Unity denotes zero phase contribution; Minimum/Maximum labels refer to precession-index extrema. C combines chronology and refitting variability but is not a calibrated combined confidence interval; B and C need not nest.}
\label{fig:effect_precision}
\end{figure}

The primary association persists across the tested specifications. Decay times from 1 to 5 kyr give $G$ 0.243--0.252 and phase maxima $318.2^\circ$--$330.3^\circ$ on identical response support; alternative pre-anchor history levels have little effect ($G$ 0.248--0.250; Figure~S9). Alternative polynomial background responses give $G$ 0.239--0.252. A rectangular history count gives $G$ 0.219 and a $332.4^\circ$ maximum; elapsed-time histories give $G$ 0.202--0.209 (Figure~S10). Segment-specific phase coefficients do not clearly improve the fit (nominal $p=0.810$). The calibrated history tests support including recent events, and two rescaled-interval diagnostics show no clear lack of fit (Text~S8).

The fit improvements from summer-solstice insolation overlap more with those from phase than do improvements from the other scalar orbital predictors (Figure~\ref{fig:orbital_comparison}; Text~S7). Added to BG, insolation gives $G$ 0.184 for the primary catalogue and 0.063 for Barker; after phase, these gains fall to 0.049 and 0.002. Conversely, phase $G$ after insolation is 0.113 and 0.050, versus 0.248 and 0.111 after BG. Neither residual phase test passes the separate nominal Holm correction (adjusted $p=0.093$ and 0.560). Eccentricity and obliquity give small additional gains after phase. Phase retains a Holm-adjusted association after obliquity ($p=0.000394$ and 0.0368), supporting the added phase terms rather than an obliquity effect.

\begin{figure}[p]
\centering
\includegraphics[width=\linewidth,height=0.68\textheight,keepaspectratio]{figures/Fig04.pdf}
\caption{Log-likelihood gain per event from orbital predictors conditional on alternative baselines for the primary catalogue (a--c) and Barker varying-threshold catalogue (d--f). BG includes event history, LR04, CO$_2$, and a primary segment intercept; Pre denotes phase sine/cosine; Orb denotes the scalar orbital variable named in each row. Columns show additions of Orb to BG, Orb after Pre, and Pre after Orb. Circles mark nominal-age $G$ (bits/event), ticks mark chronology medians, and bars span the 2.5th--97.5th percentiles of valid age fits (499 primary; 500 Barker). Dashed lines show the nominal BG + Pre versus BG comparison. Orb adds one coefficient and Pre two. Bars describe chronology sensitivity, not sampling confidence intervals; the nine additional nominal tests per catalogue use Holm correction (Text~S7).}
\label{fig:orbital_comparison}
\end{figure}

\section{Discussion}
Precession phase improves the fit to warming occurrence after accounting for background climate and recent events. This extends the Alpine interstadial--insolation comparison \citep{fohlmeister2023role} by conditioning on climate, history and exposure while assessing uncertainty. Both main catalogues favor phases near the precession-index minimum, broadly corresponding to enhanced Northern Hemisphere summer insolation. The agreement suggests probabilistic organization of warming times, although the angular offset from an index extremum cannot identify a physical response lag.

A phase-dependent rate is compatible with stochastic transitions whose likelihood changes with climate state. Internal variability can produce abrupt changes \citep{ditlevsen2005recurrence,kleppin2015stochastic}; sea-ice and ocean feedbacks or CO$_2$-dependent oscillatory regimes can alter the conditions for their occurrence \citep{dokken2013dansgaard,pedro2022dansgaard,vettoretti2022atmospheric}. Model experiments also show sensitivity to insolation configuration \citep{zhang2021direct}, without uniquely predicting the response fitted here. Within the nonpositive-history model, fitted coefficients are $-0.964$ and $-1.129$. Simulations under zero history give $p=0.026$ and $0.012$, respectively, supporting a history contribution beyond climate and phase (Text~S8). This could reflect climate-state cycling or inventory construction, rather than a specific recovery mechanism. Moreover, whole-time onset rates depend on time spent in warm states as well as the likelihood of leaving a cold state \citep{lohmann2018random}.

Occurrence and amplitude provide different tests of orbital influence. The methane--speleothem comparison of \citet{thirumalai2020methane} concerns variability envelopes that combine proxy sensitivity, event magnitude, duration and frequency. Our association therefore neither contradicts their Chinese-speleothem result nor shows that warming frequency controls methane amplitude. Comparing these distinct quantities can help determine which aspects of abrupt variability carry an orbital signal.

The orbital comparisons leave more than one physical interpretation open. Insolation depends on several orbital elements, so overlap between its fit improvement and that of phase is expected. The two-coefficient phase response and one-coefficient insolation response also impose different shapes; raw $G$ is not an equal-complexity ranking. Conditioning on LR04 and CO$_2$ may absorb pathways through which orbital forcing changes background climate. Small additional eccentricity or obliquity contributions therefore do not exclude indirect effects. Exploratory LR04--phase interactions provide no clear evidence of modulation (nominal $p=0.187$ for the primary catalogue and $p=0.0775$ for Barker; Figure~S11).

Detection and chronology assumptions limit these comparisons. Published inventories do not guarantee uniform completeness, and the regional speleothems and synthetic Barker reconstruction have different transfer functions. Additional random event deletion generally preserves the fitted phase sector (Text~S8), but this stress test cannot reconstruct unrecorded transitions or exclude phase-dependent detection bias. We condition on retained event identities and approximate pre-anchor history by zero. History decay and the nonpositive coefficient constraint are likewise working assumptions. Rescaled-interval diagnostics show no clear misfit under the two tested statistics, rather than establishing that all relevant climate or history dependence has been captured.

Chronology perturbations capture only part of the dating uncertainty. The modeled Greenland extension inherits the ss09sea framework, whose development used a marine-isotope-stage age constraint from SPECMAP \citep{southon2004radiocarbon,wolff2010millennial}; this does not mean individual NGRIP events were tuned to precession. LR04 is orbitally tuned, and Barker uses speleothem matching \citep{lisiecki2005pliocene,barker2011800}. Age perturbations cannot remove these dependencies. Nor do the MF shared factor, transferred Sofular control errors, model-extension envelope and fixed forcing ages constitute complete age-model posteriors. Finally, finite event counts permit a wide range of modulation strengths even though null simulations support the phase association. The point-estimated rate ratios should not be read as precise physical sensitivities.

\section{Conclusions}
Precession phase improves the description of warming occurrence beyond recent events, LR04 and CO$_2$ in the Greenland--speleothem catalogue, with 53 fitted response events from an inventory of 55. The association generally persists across the tested chronology and model variations, and the synthetic Greenland catalogue favors a similar phase sector. Effect strength remains imprecise, while overlap between phase and summer-insolation fits and chronology dependencies preclude identifying a unique orbital mechanism. Models of abrupt climate variability should therefore be tested against the distribution of warming times as well as event amplitudes and durations.

\section*{Abbreviations}

\begin{description}
\setlength{\itemsep}{0pt}
\setlength{\parsep}{0pt}
\item[AIC] Akaike information criterion
\item[b2k] before 2000 CE
\item[BG] background model, including climate and event-history terms
\item[BP] before present, with present defined as 1950 CE
\item[CE] Common Era
\item[GI] Greenland Interstadial
\item[GICC05] Greenland Ice-Core Chronology 2005
\item[GS] Greenland Stadial
\item[La2004] orbital solution of Laskar et al. (2004)
\item[LR] likelihood ratio
\item[LR04] benthic $\delta^{18}$O stack of Lisiecki and Raymo (2005)
\item[MC] Monte Carlo
\item[MCE] maximum counting error
\item[MF] Melchsee--Frutt
\item[MIS] Marine Isotope Stage
\item[NGRIP] North Greenland Ice Core Project
\item[Orb] scalar orbital predictor used in model comparisons
\item[Pre] precession-phase sine and cosine terms
% SPECMAP expansion: https://www.ncei.noaa.gov/products/paleoclimatology/paleoceanography
\item[SPECMAP] Spectral Mapping Project
\item[ss09sea] Greenland ice-flow-model chronology underlying the modeled extension
\item[U--Th] uranium--thorium
\end{description}

\section*{Open Research}
The Greenland event inventory follows \citet{rasmussen2014stratigraphic}; the GICC05 data are archived by \citet{rasmussen2022gicc05}. MF and Sofular source data are archived by \citet{fohlmeister2023mfData,held2024sofularData}. The synthetic Greenland event catalogue and its chronology controls are supplied by \citet{barker2011800}. Background and orbital inputs follow \citet{lisiecki2005pliocene,bereiter2015revision,laskar2004long}. Processed catalogues, observation intervals, saved model results, and reproducible figure/table scripts accompany this manuscript in the research workspace. Public code and data archive identifiers will be added before submission.

\section*{Acknowledgements and author statements}
\emph{Author names and affiliations, contributions, funding, acknowledgements, and conflict-of-interest statements await author confirmation.}

\bibliographystyle{agu}
\bibliography{reference}
\end{document}
